\scanpage{381}
\mspage{193}

\inserted{\unsure{germs of|germ of}} liftings, for \unsure{a|an} \inserted{germ of }outer action on $\pi$ already given --- which corresponds to the naive fact that the algebraic curve $U_{\deleted{\unsure{g,\nu}}}$ over $\overline{\mathbb Q}$ defined by this outer action \unsure{has} ``many point\textsuperscript{s}'' rational/$\overline{\mathbb Q}$. In fact, it would suffice, in the limit, to speak of the \deleted{\illegible{1}} properties of these more complete liftings to \inserted{a germ of }a genuine action of \deleted{\illegible{1}} $\mathbb{\Gamma}_a$ on $\pi$ (not only outer) --- from which the outer actions \unsure{are|will} \unsure{to be|be} deduced by composition with
\[
M_a \longrightarrow N_a \simeq M_a/\pi.
\]

\deleted{\illegible{1}} We thus suppose given a closed subgroup
\[
\Gamma \subset \deleted{\unsure{\mathcal G}} M_a \subset \widehat{\widehat{\mathcal G}}(\pi)
\]
such that
\[
\text{\textcircled{a}}\qquad
\Gamma \longrightarrow \mathbb{\Gamma}_a
 = M_a/\pi
\quad \text{\unsure{is|is} injective,}
\]
and \unsure{has|has} as its image an open subgroup of $\mathbb{\Gamma}_a$.

--- it being understood that two subgroups conjugate under $\pi$ --- or even, \inserted{sometimes,} under $M_a$ --- are not considered essentially distinct. We consider, at the same time as $\Gamma$, subgroups

\scanpage{382}

$\Gamma'$ that are open. We surely want, for such open subgroups
\[
\text{\textcircled{b}}\qquad \pi^{\Gamma'}=(1)
\]
If we denote by $\widetilde{\Gamma}$ the image of $\Gamma$ in $N_a$, so that we have an extension
\begin{samepage}
(1)\quad
\diagram{\begin{tikzcd}[ampersand replacement=\&,column sep=2.1em,row sep=1.1em]
1 \arrow[r,dash] \& \pi \arrow[r,dash] \& \mathcal E \arrow[r] \& \widetilde{\Gamma} \arrow[r,dash] \& 1 \\
\& \& \& \Gamma \arrow[u,"\sim"'] \&
\end{tikzcd}}
\end{samepage}
we can regard $\Gamma$ as a section of this extension, $\Gamma'$ as a partial section. The hypothesis $\pi^{\Gamma'}=(1)$ ensures the \uline{rigidity} of the category of points of $U$ with values in $\overline{\mathbb Q}$, paradigmed by that of germs of \unsure{splittings|splitting} of the extension \textcircled{(1)'}.
\margin{\unsure{!|?}}

We can also look at the sets
\[
\unsure{\widehat{\mathcal G}_a^{+\Gamma'}}
=
\deleted{\unsure{N_a^{\Gamma'}|M_a^{\Gamma'}}}
\operatorname{Centr}_{M_a}(\Gamma')\cap\widehat{\mathcal G}_a^+
\quad\text{and}
\]
\[
\widehat{\mathcal T}_a^{+\Gamma'}
=\operatorname{Centr}_{N_a}(\widetilde{\Gamma}')\cap\widehat{\mathcal T}_a^+,
\]
\deleted{and} we thus have an exact sequence
\[
(2)\qquad 1\longrightarrow\pi^{\Gamma'}\longrightarrow
(\widehat{\mathcal G}_a^+)^{\Gamma'}\longrightarrow
(\widehat{\mathcal T}_a^+)^{\Gamma'}
\]
which, taking into account the hypothesis $\pi^{\Gamma'}=(1)$,

\scanpage{383}
\mspage{194}

gives
\[
(3)\qquad (\widehat{\mathcal G}_a^+)^{\Gamma'}
\hookrightarrow (\widehat{\mathcal T}_a^+)^{\Gamma'}.
\]

The right-hand side of (3), by our hypothetical dictionary, should be \inserted{\unsure{canonically}} \unsure{the|isomorphic} group of \unsure{automorphisms} of the curve $U$ over $\unsure{K'_a|K'}$ hence be a finite group, and even \deleted{\unsure{by}} the inductive limit (for $\Gamma'$ decreasing) --- which is the group of \unsure{automorphisms} of $U$ over $\overline{\mathbb Q}$ --- is finite. Denoting by a superscript $\natural$ the corresponding germ of a group, we thus want
\[
\text{\textcircled{c}}\qquad
Z=(\widehat{\mathcal T}_a^+)^{(\Gamma^{\natural})}
\left(=\operatorname{Centr}_{\unsure{N_a|M_a}}
(\widetilde{\Gamma}^{\natural})\cap\widehat{\mathcal T}_a^+\right)
\]
is a \uline{finite} group

--- which is a counterpart to (b), and expresses that the automorphism groups of the points of $M_{g,\nu\overline{\mathbb Q}}$ over $\overline{\mathbb Q}$ are finite.

In fact, we would like $Z$ to be conjugate in $\deleted{\unsure{\widehat{\mathcal T}}_a^+}\inserted{M_a}$ to a subgroup of \unsure{to} $\mathcal T^+$ \deleted{Centr} (if we suppose that we have a \inserted{discretification} \deleted{\illegible{1}} $\pi_0$ of $\pi$, making it possible to define $\mathcal T$ --- otherwise, \unsure{we can} express this

\scanpage{384}

property by saying that there exists a \unsure{discretification} \inserted{$\pi_0$} of $\pi$, compatible with $a$, such that $Z\subset\mathcal T(\pi_0)$\ldots)

\unsure{Let us consider} now an arbitrary finite subgroup $G\subset\widehat{\mathcal T}_a^+$. [N.B. If we take \inserted{only} $G\subset M_a$, so that the image of $G$ in $\mathbb{\Gamma}_a=M_a/\widehat{\mathcal T}_a^+$ is a finite subgroup, then if it \deleted{$\mathbb{\Gamma}_a$} is true that $\mathbb{\Gamma}_a\simeq\unsure{\mathcal G|\mathfrak G}$, \margin{$\operatorname{Gal}(\overline{\mathbb Q}/\mathbb Q)\longrightarrow$} this image must be of order $1$ or $2$, and in the second case must \margin{\unsure{?}} define a $\tau\in\mathbb{\Gamma}_a$ corresponding to a well-determined \uline{\unsure{prediscretification}} (not oriented) of $\pi$. This case should be reused again in what follows\ldots]. \deleted{If we} \inserted{Let us suppose, \unsure{or}} \unsure{admit}, that $G$ is contained in $\mathcal T^+$, \margin{this is probably \uline{not} automatic (if we \unsure{do} \unsure{not|N'} suppose $N_a=\operatorname{Norm}_{\widehat{\widehat{\mathcal T}}}(\unsure{\widehat{\mathcal T}_a|\widehat{\mathcal T}_a^+})$)} for a suitable discretification in the class $a$, and that the outer action of $G$ thus obtained on a discrete $\pi_0$ of type $(g,\nu)$ \unsure{is|are} always \uline{realizable}. It is therefore also realizable \unsure{by|for} an action of $G$

\scanpage{385}
\mspage{195}

on a complex, hence algebraic, structure, which means that we can find a \deleted{\illegible{1}} \inserted{\unsure{germ|germs} of \uline{admissible} lifting}
\[
\mathbb{\Gamma}_a\longrightarrow N_a
\]
which centralizes $G$. Let us consider on the other hand
\[
(4)\qquad
\underset{\inserted{\widehat N_a}}{\operatorname{Centr}_{N_a}}(G)
\longrightarrow\mathbb{\Gamma}_a
\]
whose kernel is $\operatorname{Centr}_{\widehat{\mathcal T}_a^+}(G)=\widehat{\mathcal T}_a^{+G}$.
If we \unsure{placed|placing} ourselves in the discrete context \deleted{\unsure{for}} (with a \inserted{discretification} $\pi_0\subset\pi$ \unsure{invariant|invariant} under $G$) we would have, by the standard topological conjectures, that $\unsure{\widehat{\mathcal T}_a^{+G}|\widehat{\widehat{\mathcal T}}_a^{+G}}$ is the group of automorphisms, in the isotopic category, of the action of $G$ on a surface $U$ of type $(g,\nu)$, \unsure{described} by the given outer operation of $G$ on $\pi_0$. This group has no reason to be finite --- if it were, the \unsure{homomorphism} (4) would have finite kernel, hence would be an \unsure{isomorphism} \unsure{of the|of} kernels of groups, and there would be a

\scanpage{386}

unique germ of lifting $\mathbb{\Gamma}_a^{\natural}\longrightarrow M_a$ which would centralize $G$ --- which would mean that there is \deleted{only one} only one way to realize the \inserted{topological} operation of $G$ on $U$ by a \deleted{algebraic} complex analytic operation, hence algebraic --- but this is surely not the case, the \unsure{set} of fixed \unsure{points} of $G$ operating on Teichmüller space $\widetilde M_{g,\nu}$ is not reduced to a point --- it is \inserted{(in the topological context)} a \unsure{schematic|arithmetic} multiplicity which can be \deleted{\illegible{2}} of arbitrary $\unsure{\mathbb R|\mathbb C}$ dimension --- It is of dimension $>0$ in any case, if the quotient \deleted{\illegible{1}} $U/G$ is not of genus $0$.

To keep in mind in any case, as a plausible property :

(d) For $\forall$ finite subgroup \inserted{$G$} of $\widehat{\mathcal T}_a^+$ (or at least if $G\subset\mathcal T^+$, when we have a \inserted{\unsure{discretification}} $\pi_0\subset\pi$ \inserted{\unsure{in $a$}}), looking at $\operatorname{Centr}_{\unsure{\widehat{\widehat{\mathcal T}}|\widehat{\mathcal T}}}(G)=\unsure{;}Z(G)$, $N_a\cap Z(G)\inserted{\longrightarrow\mathbb{\Gamma}_a}$ \deleted{is open in $N_a$} \inserted{has an open image (and there are \unsure{even} some}

\scanpage{387}
\mspage{196}

germs of admissible liftings $\mathbb{\Gamma}_a\to Z(G)$)

This implies a nontrivial property for
the $g\in\unsure{N_a|M_a}$ \illegible{1} as elements of $\widehat{\widehat{\mathcal T}}(\pi)$,
or better of $\inserted{N'=}\operatorname{Norm}_{\widehat{\widehat{\mathcal T}}}(\widehat{\mathcal T})$ :
\deleted{$\exists n\in\mathbb N^{*}$ such that $g^n$ \unsure{,|;}}
\margin{\unsure{namely} : $\exists n\in\mathbb N^{*}$ such that $g^n$ is congruent modulo $\widehat{\mathcal T}^{+}_a$ to an element of $Z(G)$\ldots}
Let us indeed consider the image \inserted{$\unsure{Z'_G|Z'(G)}$} of
$Z(G)\cap N'$ in $\Gamma'=N'/\widehat{\mathcal T}$, \deleted{\unsure{then is|then let}}
\inserted{where obviously $\mathbb{\Gamma}_a\subset\Gamma'$, and \deleted{\unsure{the property}} the image of $Z(G)\cap N_a$ in $\mathbb{\Gamma}_a$ is none other than}
\deleted{\illegible{1} let $\mathbb{\Gamma}_a$ be that of $\Gamma$}
$Z'_G\cap\mathbb{\Gamma}_a$. Saying that it is \inserted{\unsure{open|open} i.e. \uline{of index}} \unsure{finite|finite},
therefore implies that $\forall g\in\mathbb{\Gamma}_a$, \deleted{\illegible{1}} $\exists n\in\mathbb N^{*}$
such that $g^n\in Z'_G$. \deleted{We \unsure{To tell the truth}}

Let us note that in $\mathcal T(\pi_{g,\nu})$ there are
only a finite number of conjugacy classes
of finite subgroups, so if we
look at their centralizers \inserted{$Z_G$} in
$\widehat{\widehat{\mathcal T}}(\pi_{g,\nu})$, and their images in
$\Gamma'_{g,\nu}=\operatorname{Norm}_{\widehat{\widehat{\mathcal T}}_{g,\nu}}(\widehat{\mathcal T}^{+}_{g,\nu})/\widehat{\mathcal T}^{+}_{g,\nu}$, we
find only a finite number of subgroups
of $\Gamma'_{g,\nu}$, let $Z'_{g,\nu}$ be their intersection.
We thus see that \unsure{this|the} subgroup $\mathbb{\Gamma}_{g,\nu}$
of $\Gamma'_{g,\nu}$ is such that $\mathbb{\Gamma}_{g,\nu}\cap Z'_{g,\nu}$ is
open in $\mathbb{\Gamma}_{g,\nu}$

\scanpage{388}

Let us now pass \unsure{to} finite
subgroups of $\widehat{\widehat{\mathcal G}}(\pi)$ itself --- or at least
of $M_a$ --- or, what comes to almost the
same thing, of $\widehat{\mathcal G}^{+}_a$. We shall, as
earlier, restrict ourselves to $G\subset\mathcal G(\pi_0)$,
for a suitable \unsure{discretification} $\in a$,
for otherwise there would be hardly anything to
say. We know that then $G$ is
cyclic --- \unsure{it|and} corresponds to an action
of $G$ on a topological $U$, \uline{with}
\uline{fixed point}. We can realize it in
a complex way, which we \unsure{will express|express} by
saying that there exists an admissible
lifting $\mathbb{\Gamma}_a^{\natural}\longrightarrow\deleted{\unsure{\mathcal G}} M_a$, which
centralizes $G$. We know moreover, if
$G\ne1$, that \deleted{$\unsure{\pi}$}
\[
(5)\qquad \pi^G=(1)
\]
(or rather, we know it in the discrete case ---
we admit it in the profinite context) ---

\scanpage{389}
\mspage{197}

which we can \unsure{again} interpret as a rigidity property --- thus
\[
\operatorname{Centr}_{\unsure{\widehat{\widehat{\mathcal G}}}}(G)
=\widehat{\widehat{\mathcal G}}^{G}
\longrightarrow\widehat{\widehat{\mathcal T}}^{G}
\]
is \uline{injective}, so if we have a $\unsure{\Gamma''|\widetilde{\Gamma}}\subset\unsure{N_a|M_a}$ which is in the image, then there exists a unique lifting $\unsure{\Gamma''|\widetilde{\Gamma}}\longrightarrow\widehat{\widehat{\mathcal G}}$ which centralizes $G$ i.e. such that it is $\Gamma\longrightarrow\widehat{\widehat{\mathcal G}}^{G}$.

But in fact what we shall know, for a given $\Gamma\subset M_a$ (\unsure{defining} an algebraic curve via its outer action on a $\pi_1$) is that $\Gamma\subset M_a^G$ (i.e. the outer action commutes with a given outer action of $G$ \inserted{i.e. $G$ operates on the algebraic curve}) --- and we would nevertheless like to conclude from this that when we have lifted $G\longrightarrow\widehat{\mathcal T}^{+}$ \unsure{to|or} $G\longrightarrow\widehat{\mathcal G}^{+}$ (i.e. when we have given ourselves a fixed point of the action of $G$ on $U$): \unsure{then|therefore} the action of $\Gamma^{\natural}$ lifts \inserted{(i.e. $\Gamma^{\natural}\longrightarrow M_a^G$)} \unsure{automatically|automatic}, \inserted{to $\Gamma^{\natural}\longrightarrow M_a$} in such a way \unsure{as to} \unsure{commute|commuting|commuting},\ldots . Is this a property \unsure{of the choice of the|of the choices of the} $M_a$, or of that of the lifting $\Gamma^{\natural}\unsure{\longrightarrow|\hookrightarrow}N_a$, or of $G$??

\scanpage{390}

In the discrete setting, unless I am mistaken, if we have a faithful action \inserted{\unsure{direct|discrete},} of finite $G$ on $\pi_0$, then \deleted{Centr}
\[
\mathcal G(\pi_0)^G\longrightarrow\mathcal T(\pi_0)^G
\]
(which is injective, by $\pi_0^G=1$) is also surjective. No, there is a mistake --- that \unsure{would mean} as it stands that if $G$ operates faithfully on a topological surface $U$ of type $g,\nu$, with a fixed point $P$, then the automorphisms which commute with $G$ \uline{fix} $P$ (instead of merely permuting the fixed points among themselves\ldots). So we must not expect $\forall$ \unsure{element} of $\unsure{N_a^G|M_a^G}$, nor even of $\widehat{\mathcal T}_a^G$, to lift to $M_a^G$ --- but rather this: for $G\unsure{\in|\subset}\mathcal T_a(\pi_0)$ a given finite \unsure{subgroup}, the $\pi$-conjugacy classes of liftings to $\widehat{\mathcal G}$, \deleted{\unsure{to} $\mathcal G$} which are interpreted as fixed points of an action of $G$ on some $U$, form a finite set, on which $\widehat{\widehat{\mathcal T}}(\pi)^G$ operates in a natural way, and \unsure{the|the} \unsure{stabilizers|stabilizer} of a point $P$ i.e. of a conjugacy class of \unsure{lifting\textsuperscript{s}|liftings}
